\documentclass[12pt,a4paper]{article}

\usepackage[utf8]{inputenc}
\usepackage[vietnamese]{babel}
\usepackage{amsmath,amssymb,amsthm}
\usepackage{geometry}
\usepackage{enumitem}
\usepackage[most]{tcolorbox}
\usepackage{hyperref}
\usepackage{fancyhdr}
\usepackage{tikz}
\usepackage{eso-pic}
\usepackage{xcolor}
\usepackage{array}

\geometry{a4paper, margin=2cm, bottom=2.5cm}
\hypersetup{colorlinks=true, linkcolor=blue!70!black, urlcolor=blue}
\setlist[itemize]{topsep=4pt, itemsep=2pt}
\setlist[enumerate]{topsep=4pt, itemsep=2pt}

\AddToShipoutPictureFG{%
  \AtPageCenter{%
    \begin{tikzpicture}[remember picture, overlay]
      \node[rotate=45, scale=1, opacity=0.06]{\fontsize{60}{72}\selectfont\bfseries\sffamily Đặng Tấn Quí};
    \end{tikzpicture}%
  }%
}

\pagestyle{fancy}
\fancyhf{}
\fancyhead[L]{\small\textit{Lý thuyết Đồ thị \& Thuật toán}}
\fancyhead[R]{\small\textit{Đặng Tấn Quí}}
\fancyfoot[C]{\thepage}
\fancyfoot[L]{\footnotesize\textcopyright\ Đặng Tấn Quí}
\fancyfoot[R]{\footnotesize SĐT: 0377\,122\,210}
\renewcommand{\headrulewidth}{0.4pt}
\renewcommand{\footrulewidth}{0.4pt}
\fancypagestyle{plain}{\fancyhf{}\fancyfoot[C]{\thepage}\fancyfoot[L]{\footnotesize\textcopyright\ Đặng Tấn Quí}\fancyfoot[R]{\footnotesize SĐT: 0377\,122\,210}\renewcommand{\headrulewidth}{0pt}\renewcommand{\footrulewidth}{0.4pt}}

\newtcolorbox{dinhnghia}[1][]{enhanced, breakable, colback=blue!3!white, colframe=blue!60!black, fonttitle=\bfseries, title={Định nghĩa}, #1}
\newtcolorbox{dinhly}[1][]{enhanced, breakable, colback=green!3!white, colframe=green!50!black, fonttitle=\bfseries, title={Định lý}, #1}
\newtcolorbox{tinhchat}[1][]{enhanced, breakable, colback=yellow!5!white, colframe=orange!50!black, fonttitle=\bfseries, title={Tính chất}, #1}
\newtcolorbox{phuongphap}[1][]{enhanced, breakable, colback=gray!5!white, colframe=gray!50!black, fonttitle=\bfseries, title={Phương pháp}, #1}
\newtcolorbox{luuy}[1][]{enhanced, breakable, colback=red!3!white, colframe=red!50!black, fonttitle=\bfseries, title={Lưu ý}, #1}
\newtcolorbox{congthuc}[1][]{enhanced, breakable, colback=cyan!3!white, colframe=cyan!50!black, fonttitle=\bfseries, title={Công thức}, #1}
\newtcolorbox{vidu}[1][]{enhanced, breakable, colback=violet!3!white, colframe=violet!50!black, fonttitle=\bfseries, title={Ví dụ}, #1}

\begin{document}

\thispagestyle{empty}
\begin{tikzpicture}[remember picture, overlay]
  \fill[blue!80!black] (current page.south west) rectangle (current page.north east);
  \fill[blue!60!cyan, opacity=0.8] (current page.south west) -- (current page.north west) -- ([xshift=-3cm]current page.north east) -- (current page.south east) -- cycle;
  \foreach \r/\o in {6/0.05, 4.5/0.08, 3/0.12} {\fill[white, opacity=\o] ([xshift=2cm, yshift=-2cm]current page.north east) circle (\r cm);}
  \draw[white, line width=2pt] ([yshift=-5cm]current page.north west) ++(2cm,0) -- ++(17cm,0);
  \draw[white, line width=2pt] ([yshift=-16cm]current page.north west) ++(2cm,0) -- ++(17cm,0);
  \node[anchor=west, white, font=\fontsize{24}{30}\bfseries\selectfont] at ([xshift=3cm, yshift=-7.5cm]current page.north west) {PHẦN 12: LÝ THUYẾT};
  \node[anchor=west, white, font=\fontsize{20}{26}\bfseries\selectfont] at ([xshift=3cm, yshift=-9.2cm]current page.north west) {ĐỒ THỊ \& THUẬT TOÁN};
  \node[anchor=west, white, font=\fontsize{14}{18}\selectfont] at ([xshift=3cm, yshift=-10.7cm]current page.north west) {(Decision Mathematics)};
  \node[anchor=west, white, font=\fontsize{16}{20}\selectfont] at ([xshift=3cm, yshift=-12cm]current page.north west) {Tài liệu luyện thi du học};
  \node[anchor=west, white, font=\Large\bfseries] at ([xshift=3cm, yshift=-13.5cm]current page.north west) {Biên soạn:};
  \node[anchor=west, yellow!80!white, font=\fontsize{28}{34}\bfseries\selectfont] at ([xshift=3cm, yshift=-15cm]current page.north west) {ĐẶNG TẤN QUÍ};
  \node[anchor=west, white!90, font=\large] at ([xshift=3cm, yshift=-18cm]current page.north west) {SĐT: 0377\,122\,210};
  \node[anchor=south, white!80, font=\small] at ([yshift=1.5cm]current page.south) {Tài liệu có bản quyền --- Cấm sao chép khi chưa được phép};
  \node[anchor=south east, white, font=\Large\bfseries] at ([xshift=-2cm, yshift=3cm]current page.south east) {\the\year};
\end{tikzpicture}
\newpage

\thispagestyle{empty}
\vspace*{5cm}
\begin{center}
  {\Large\bfseries PHẦN 12: LÝ THUYẾT ĐỒ THỊ \& THUẬT TOÁN}\\[8pt]
  {\large (Decision Mathematics)}\\[4pt]
  {\large Tài liệu luyện thi du học}
\end{center}
\vspace{2cm}
\begin{tcolorbox}[enhanced, colback=red!3!white, colframe=red!60!black, fonttitle=\bfseries, title={Thông tin bản quyền}, boxrule=1.5pt]
  \textbf{Tác giả:} Đặng Tấn Quí \quad \textbf{Liên hệ:} 0377\,122\,210 \quad \textbf{Năm:} \the\year \\[6pt]
  \textbf{Bản quyền \textcopyright\ \the\year\ Đặng Tấn Quí. Bảo lưu mọi quyền.}
\end{tcolorbox}
\vfill\begin{center}\small\textit{Tài liệu lưu hành nội bộ}\end{center}
\newpage
\tableofcontents
\newpage

%% ================================================================
%%  PHẦN 12: LÝ THUYẾT ĐỒ THỊ & THUẬT TOÁN
%% ================================================================
\section{Lý thuyết Đồ thị \& Thuật toán (Decision Mathematics)}

%% ----------------------------------------------------------------
%%  12.1 Thuật toán tìm đường ngắn nhất: Dijkstra, Prim, Kruskal
%% ----------------------------------------------------------------
\subsection{Thuật toán tìm đường ngắn nhất và cây khung tối tiểu}

\subsubsection{Khái niệm cơ bản về đồ thị}

\begin{dinhnghia}[title={Đồ thị (Graph)}]
  \textbf{Đồ thị} $G = (V, E)$ gồm:
  \begin{itemize}
    \item $V$: tập đỉnh (vertices/nodes).
    \item $E$: tập cạnh (edges), mỗi cạnh nối hai đỉnh.
  \end{itemize}
  \textbf{Phân loại:}
  \begin{itemize}
    \item \textbf{Đồ thị vô hướng:} cạnh không có hướng.
    \item \textbf{Đồ thị có hướng (digraph):} mỗi cạnh có hướng.
    \item \textbf{Đồ thị có trọng số:} mỗi cạnh gán một trọng số (weight) thể hiện khoảng cách, chi phí.
    \item \textbf{Đồ thị liên thông:} tồn tại đường đi giữa mọi cặp đỉnh.
    \item \textbf{Cây (Tree):} đồ thị liên thông không có chu trình, $|E| = |V| - 1$.
  \end{itemize}
\end{dinhnghia}

\subsubsection{Thuật toán Dijkstra (Đường đi ngắn nhất)}

\begin{phuongphap}[title={Thuật toán Dijkstra}]
  Tìm đường đi ngắn nhất từ đỉnh nguồn $s$ đến tất cả các đỉnh khác trong đồ thị \textbf{có trọng số không âm}.

  \textbf{Thuật toán:}
  \begin{enumerate}
    \item \textbf{Khởi tạo:} Gán $d(s) = 0$, $d(v) = \infty$ với mọi $v \ne s$. Tập $S = \emptyset$ (đỉnh đã xử lý), $Q = V$ (đỉnh chưa xử lý).
    \item \textbf{Lặp:} Trong khi $Q \ne \emptyset$:
      \begin{enumerate}
        \item Chọn đỉnh $u \in Q$ có $d(u)$ nhỏ nhất.
        \item Chuyển $u$ từ $Q$ sang $S$ (đánh dấu đã xử lý).
        \item Với mỗi láng giềng $v$ của $u$ chưa xử lý:
          \[
            \text{Nếu } d(u) + w(u,v) < d(v): \quad d(v) \leftarrow d(u) + w(u,v),\quad \text{prev}(v) \leftarrow u.
          \]
      \end{enumerate}
    \item \textbf{Kết quả:} $d(v)$ là độ dài đường đi ngắn nhất từ $s$ đến $v$.
  \end{enumerate}

  \textbf{Độ phức tạp:} $O(V^2)$ (đơn giản) hoặc $O((V+E)\log V)$ (dùng hàng đợi ưu tiên).
\end{phuongphap}

\begin{vidu}
  Tìm đường đi ngắn nhất từ $A$ đến tất cả đỉnh trong đồ thị:

  \begin{center}
    \begin{tikzpicture}[every node/.style={circle, draw, minimum size=6mm}]
      \node (A) at (0,2) {A};
      \node (B) at (2,3) {B};
      \node (C) at (2,1) {C};
      \node (D) at (4,3) {D};
      \node (E) at (4,1) {E};
      \draw (A) -- node[above, draw=none] {4} (B);
      \draw (A) -- node[below, draw=none] {2} (C);
      \draw (B) -- node[above, draw=none] {3} (D);
      \draw (B) -- node[left, draw=none] {1} (C);
      \draw (C) -- node[below, draw=none] {5} (E);
      \draw (D) -- node[right, draw=none] {2} (E);
    \end{tikzpicture}
  \end{center}

  \medskip\textbf{Giải:}
  \begin{itemize}
    \item Khởi tạo: $d(A)=0$, $d(B)=d(C)=d(D)=d(E)=\infty$.
    \item Chọn $A$ ($d=0$): Cập nhật $d(B)=4$, $d(C)=2$.
    \item Chọn $C$ ($d=2$): Cập nhật $d(B)=\min(4,2+1)=3$, $d(E)=7$.
    \item Chọn $B$ ($d=3$): Cập nhật $d(D)=\min(\infty,3+3)=6$.
    \item Chọn $D$ ($d=6$): Cập nhật $d(E)=\min(7,6+2)=7$ (không đổi).
    \item Chọn $E$ ($d=7$).
  \end{itemize}
  Kết quả: $d(A)=0$, $d(C)=2$, $d(B)=3$, $d(D)=6$, $d(E)=7$.
\end{vidu}

\subsubsection{Cây khung nhỏ nhất (Minimum Spanning Tree)}

\begin{dinhnghia}[title={Cây khung (Spanning Tree)}]
  Cây khung của đồ thị $G = (V,E)$ là cây (đồ thị liên thông, không chu trình) gồm tất cả $|V|$ đỉnh và $|V|-1$ cạnh của $G$.

  \textbf{Cây khung nhỏ nhất (MST):} cây khung có tổng trọng số nhỏ nhất.
\end{dinhnghia}

\begin{phuongphap}[title={Thuật toán Prim (Xây dựng MST từ một đỉnh)}]
  \begin{enumerate}
    \item Bắt đầu từ đỉnh tùy ý, thêm vào tập $T$ (cây đang xây dựng).
    \item Lặp cho đến khi $T$ có $|V|$ đỉnh:
      \begin{itemize}
        \item Tìm cạnh trọng số nhỏ nhất nối đỉnh trong $T$ với đỉnh ngoài $T$.
        \item Thêm cạnh đó (và đỉnh mới) vào $T$.
      \end{itemize}
  \end{enumerate}
  \textbf{Nguyên lý:} Luôn chọn cạnh có trọng số nhỏ nhất ra ngoài cây hiện tại.
\end{phuongphap}

\begin{phuongphap}[title={Thuật toán Kruskal (Xây dựng MST theo cạnh)}]
  \begin{enumerate}
    \item Sắp xếp tất cả cạnh theo trọng số tăng dần.
    \item Lặp qua từng cạnh theo thứ tự:
      \begin{itemize}
        \item Nếu thêm cạnh này vào cây \textbf{không tạo ra chu trình}: thêm vào.
        \item Ngược lại: bỏ qua cạnh này.
      \end{itemize}
    \item Dừng khi cây có $|V|-1$ cạnh.
  \end{enumerate}
  \textbf{Dùng Union-Find (Disjoint Set)} để kiểm tra chu trình hiệu quả.

  \textbf{Độ phức tạp:} $O(E \log E)$.
\end{phuongphap}

\begin{vidu}
  Tìm MST của đồ thị với các cạnh: $AB=4$, $AC=2$, $BC=1$, $BD=3$, $CE=5$, $DE=2$.

  \medskip\textbf{Dùng Kruskal:}
  Sắp xếp cạnh: $BC(1)$, $AC(2)$, $DE(2)$, $BD(3)$, $AB(4)$, $CE(5)$.
  \begin{enumerate}
    \item $BC(1)$: Thêm (không chu trình). Cây: $\{BC\}$.
    \item $AC(2)$: Thêm. Cây: $\{BC, AC\}$.
    \item $DE(2)$: Thêm. Cây: $\{BC, AC, DE\}$.
    \item $BD(3)$: Thêm (nối hai thành phần $\{A,B,C\}$ và $\{D,E\}$). Cây: $\{BC, AC, DE, BD\}$.
    \item Cây có $5-1=4$ cạnh. Dừng.
  \end{enumerate}
  MST có tổng trọng số: $1+2+2+3 = 8$.
\end{vidu}

%% ----------------------------------------------------------------
%%  12.2 Quy hoạch tuyến tính
%% ----------------------------------------------------------------
\subsection{Quy hoạch tuyến tính (Linear Programming)}

\begin{dinhnghia}[title={Bài toán quy hoạch tuyến tính}]
  \textbf{Quy hoạch tuyến tính (Linear Programming -- LP)} là bài toán tìm cực trị (lớn nhất hoặc nhỏ nhất) của một \textbf{hàm mục tiêu tuyến tính} đối với các \textbf{ràng buộc tuyến tính}.

  \textbf{Dạng chuẩn:}
  \[
    \text{Tối ưu: } z = c_1 x_1 + c_2 x_2 + \cdots + c_n x_n
  \]
  \[
    \text{Ràng buộc: } \begin{cases} a_{11}x_1 + \cdots + a_{1n}x_n \le b_1 \\ \vdots \\ a_{m1}x_1 + \cdots + a_{mn}x_n \le b_m \\ x_1, x_2, \ldots, x_n \ge 0 \end{cases}
  \]
\end{dinhnghia}

\begin{phuongphap}[title={Phương pháp đồ thị (cho 2 biến)}]
  Dùng khi chỉ có hai biến $x$ và $y$:
  \begin{enumerate}
    \item \textbf{Xác định miền khả thi:} Vẽ tất cả các bất đẳng thức ràng buộc trên hệ trục $Oxy$; miền nghiệm là giao của tất cả nửa mặt phẳng.
    \item \textbf{Tìm các đỉnh (vertex) của miền khả thi:} Giải từng cặp phương trình ràng buộc để tìm các đỉnh góc.
    \item \textbf{Áp dụng Định lý đỉnh:} Giá trị cực trị của hàm mục tiêu tuyến tính đạt được tại một \textbf{đỉnh} của miền khả thi.
    \item Tính $z$ tại mỗi đỉnh, chọn giá trị lớn nhất (hoặc nhỏ nhất).
  \end{enumerate}
\end{phuongphap}

\begin{dinhly}[title={Định lý điểm cực (Corner Point Theorem)}]
  Nếu bài toán LP có nghiệm tối ưu, thì nghiệm tối ưu đạt được tại \textbf{ít nhất một đỉnh} (điểm góc) của miền khả thi lồi.
\end{dinhly}

\begin{vidu}
  Một công ty sản xuất hai loại sản phẩm $A$ và $B$. Mỗi $A$ lãi $5$ đơn vị, mỗi $B$ lãi $4$ đơn vị. Ràng buộc:
  \begin{itemize}
    \item Nguyên liệu: $2A + B \le 14$.
    \item Thời gian: $A + 2B \le 14$.
    \item Không âm: $A \ge 0$, $B \ge 0$.
  \end{itemize}
  Tối đa hóa lợi nhuận $z = 5A + 4B$.

  \medskip\textbf{Giải:}
  Tìm các đỉnh của miền khả thi:
  \begin{itemize}
    \item Giao $A=0$ và $B=0$: $(0,0)$, $z=0$.
    \item Giao $2A+B=14$ và $B=0$: $A=7$, $(7,0)$, $z=35$.
    \item Giao $2A+B=14$ và $A+2B=14$: Giải hệ: $3A=14 \Rightarrow A=14/3$, ... giải ra: $A=\dfrac{14}{3}$... thực ra:
      $\begin{cases}2A+B=14\\A+2B=14\end{cases} \Rightarrow$ trừ: $A-B=0 \Rightarrow A=B$. Thế vào: $3A=14 \Rightarrow A=B=\dfrac{14}{3}$. $z = 5\cdot\dfrac{14}{3}+4\cdot\dfrac{14}{3} = \dfrac{9\cdot14}{3} = 42$.
    \item Giao $A+2B=14$ và $A=0$: $B=7$, $(0,7)$, $z=28$.
  \end{itemize}

  So sánh: $z_{\max} = 42$ tại $A = B = \dfrac{14}{3} \approx 4.67$.

  \textbf{Nếu $A$, $B$ phải là số nguyên}: kiểm tra $(4,5)$: $z=40$; $(5,4)$: $z=41$; $(4,4)$ cần kiểm tra ràng buộc: $2(5)+(4)=14\le14$ ✓, $5+2(4)=13\le14$ ✓; $z=41$. Vậy $z_{\max}=41$ tại $(5,4)$ (hoặc kiểm tra thêm).
\end{vidu}

\begin{vidu}
  \textbf{Bài toán pha trộn (Mixture Problem):}
  Một nhà dinh dưỡng cần pha trộn hai loại thức ăn $F_1$ và $F_2$ sao cho:
  \begin{itemize}
    \item Protein tối thiểu: $3x + 5y \ge 15$.
    \item Carb tối thiểu: $x + y \ge 4$.
    \item Chi phí tối thiểu: $z = 2x + 3y$ ($x, y \ge 0$).
  \end{itemize}

  \medskip\textbf{Giải:}
  Đây là bài toán \textbf{tối thiểu hóa} với ràng buộc $\ge$ (miền khả thi không giới nội).

  Đỉnh miền khả thi:
  \begin{itemize}
    \item Giao $3x+5y=15$ và $y=0$: $(5, 0)$, $z=10$.
    \item Giao $3x+5y=15$ và $x+y=4$: $3x+5(4-x)=15 \Rightarrow -2x=-5 \Rightarrow x=2.5$, $y=1.5$. $z=5+4.5=9.5$.
    \item Giao $x+y=4$ và $x=0$: $(0,4)$, $z=12$.
  \end{itemize}

  $z_{\min} = 9.5$ tại $(2.5,\, 1.5)$.
\end{vidu}

\begin{luuy}
  \begin{itemize}
    \item Nếu miền khả thi \textbf{rỗng}: bài toán vô nghiệm.
    \item Nếu miền khả thi \textbf{không bị giới hạn} (unbounded) và tối đa hóa: có thể không có cực đại (cần kiểm tra).
    \item Bài toán \textbf{LP nguyên} (Integer LP -- $x, y \in \mathbb{Z}$): phức tạp hơn, nghiệm tối ưu không nhất thiết tại đỉnh của LP liên tục.
  \end{itemize}
\end{luuy}

\medskip\noindent\textbf{Các dạng bài tập tổng hợp:}
\begin{enumerate}[label=\textbf{Dạng \arabic*:}, leftmargin=*, align=left]
  \item Tìm đường đi ngắn nhất bằng thuật toán Dijkstra (cho bảng hoặc đồ thị).
  \item Tìm MST bằng thuật toán Prim (từ đỉnh bất kỳ).
  \item Tìm MST bằng thuật toán Kruskal (theo cạnh).
  \item So sánh Prim và Kruskal trên cùng một đồ thị.
  \item Tối đa hóa hàm mục tiêu LP hai biến (tìm đỉnh và tính $z$).
  \item Tối thiểu hóa chi phí LP hai biến.
  \item Biện luận khi thay đổi hệ số hàm mục tiêu hoặc ràng buộc.
  \item LP nguyên (Integer Programming) -- kiểm tra các điểm nguyên gần đỉnh tối ưu.
\end{enumerate}

\end{document}
